% ID: SEL-D0027-OB-E02
% ESTADO: cotejada. Segunda pasada visual completa y borrador compilado revisados.
% Fuente: Selectividad/matematicas_aplicadas/suplemente_septiembre_EXAMEN 6-ANDALUCÍA.pdf
% SHA-256: 601fd1292e7d6b21ceecde2856a283e0e929d02d4e314d97fd7062bda846650a
% Página: 2; opción B; ejercicio original 2.

\begin{enumerate}[label=\alph*)]
\item \textbf{(1.5 puntos)} Calcule los valores de los parámetros $a$ y $b$ para que la gráfica de la función $f(x)=x^3+ax^2+b$ presente un extremo relativo en el punto $(2,6)$.
\item \textbf{(1 punto)} Para $a=1$ y $b=1$, halle la ecuación de la recta tangente a la gráfica de esa función en el punto de abscisa $x=1$.
\end{enumerate}
